\documentclass[conference]{IEEEtran}
\usepackage{times}
\usepackage[numbers]{natbib}
\usepackage{graphicx}
\usepackage{booktabs}
\usepackage[bookmarks=true]{hyperref}

\begin{document}
\small

\title{Rebuttal for CoRL Submission \#416\vspace{-0.9cm}}
\maketitle
\thispagestyle{empty}
\vspace{-0.15cm}

{\small We thank the AC and all reviewers for constructive comments.
We clarify review questions as follows (R1=BKKY, R2=dWbG, R3=4zfY). }

\section*{1. Additional Baselines \& Downstream Value}

% \textbf{Q1 (AC, R2, R3): Why no comparison with ActNeRF/PB-NBV, or why is a direct comparison infeasible?} Because our contribution is the planner, we
% add a controlled algorithm-level comparison rather than reproduce the prior
% methods' different robots and reconstruction backends. We adapt their core NBV
% criteria: ActNeRF~\cite{actnerf} scores ensemble RGB-rendering variance, while
% PB-NBV~\cite{pbnbv} scores projected occupied/frontier ellipsoids. Fixed,
% Pose-Novelty, both adapted baselines, and Ray-GPIS share the same RGB-D input,
% object poses, 256 candidate views, three rotation primitives, NBV-to-action
% mapper, fusion, and evaluator; only the planner changes. The benchmark contains
% eight objects, 15 shared initial poses, and 120 paired episodes, each with one
% initial image and five planner-selected images. Compared with ActNeRF,
% Ray-GPIS reduces planning time by 88.3\% while increasing score--gain
% correlation by 0.310 (0.162 to 0.472) and reducing oracle regret by 72.2\%
% (0.0317 to 0.0088). This shows that its score is both more efficient to compute
% and better aligned with the realized geometric gain of executable actions. PB-NBV's freely selected camera views transfer poorly after mapping to
% three coarse in-hand primitives, as reflected by its negative score--gain
% correlation; this explains the result under our action constraint rather than
% implying that PB-NBV is generally inferior.

% \textbf{Q2 (AC, R1, R2): Would a simpler pose-conditioned reorientation
% achieve similar performance?} We add a \textbf{Pose-Novelty} baseline that uses
% only tracked pose: for each of the same three primitives, it predicts the next
% object-frame viewing direction and maximizes its minimum angular distance to
% the pose history. Under GT pose and exact registration, Pose-Novelty achieves
% a slightly higher F-AUC (0.9090 vs. 0.8987), because a recorded direction is nearly
% equivalent to acquired geometry. This equivalence can fail in-hand: occlusion
% and pose/depth misregistration may advance pose history while local depth
% never enters the reconstruction. We therefore retain the ideal result and make
% only a conditional robustness claim. In the targeted visited-but-unfused test,
% Pose-Novelty marks the missing direction as covered, whereas Ray-GPIS retains
% the reconstruction gap; missing-patch Recall@5 rises from 0.2827 to 0.5066
% (+79.2\%). Thus, pose coverage is effective when sensing is ideal, while
% reconstruction conditioning is more robust when nominal viewpoint coverage
% and acquired geometry diverge.

% \section*{1. Positioning, Baselines & Downstream Value}

\textbf{Q1 (AC, R2, R3): Do we compare with ActNeRF/PB-NBV, and how does AURORA relate to Scalable Real2Sim?}
We now directly compare with adapted ActNeRF and PB-NBV, but do not benchmark
Scalable Real2Sim~\cite{pfaff2025_scalable_real2sim} because it has no adaptive
view planner. Because our contribution is the planner, we isolate it in a
controlled algorithm-level comparison instead of reproducing different robots
and reconstruction backends. We adapt the NBV criteria: ActNeRF~\cite{actnerf}
uses ensemble RGB-rendering variance, while PB-NBV~\cite{pbnbv} scores projected
occupied ellipsoids. Fixed, Pose-Novelty, both adapted baselines, and Ray-GPIS
share the same RGB-D observations, poses, 256 candidate views, three rotation
primitives, NBV-to-action mapping, fusion, and evaluator; only the planner
changes. We evaluate 120 paired episodes over eight objects and 15 initial
poses, with one initial and five selected observations.

Compared with ActNeRF, Ray-GPIS reduces planning time by 88.3\%, raises score--gain correlation from 0.162 to 0.472, and reduces oracle regret by 72.2\% (0.0317 to 0.0088), indicating faster and better action ranking. PB-NBV yields negative correlation after its free camera views are mapped to three coarse in-hand primitives; this reflects our action constraint rather than general inferiority.

Scalable Real2Sim scans objects using prescribed rotations and re-grasps; its
information-optimized trajectory identifies inertial parameters rather than
selecting camera views. Its acquisition strategy is therefore closer to the
pose-conditioned reorientation evaluated in Q2 than to AURORA's
reconstruction-conditioned active planning.

\textbf{Q2 (AC, R1, R2): Would a simpler pose-conditioned reorientation
achieve similar performance?} Yes under ideal pose and registration, but not
when pose coverage and fused geometry diverge. We add a \textbf{Pose-Novelty} baseline that
selects, among the same three primitives, the predicted next viewing direction
farthest from pose history. Under GT pose and exact registration, it slightly
exceeds Ray-GPIS in F-AUC (0.9090 vs.\ 0.8987), as expected when visited
viewpoints nearly imply acquired geometry. This comparison concerns high-level
target selection: Ray-GPIS determines \emph{which} orientation should be
explored, whereas Hora-like controllers~\cite{hora} address \emph{how} to
execute the requested rotation; the two are complementary rather than
competing alternatives.

This equivalence breaks under occlusion or pose/depth misregistration: pose
history may advance while local depth is not fused. Pose-Novelty then marks the
direction as covered, whereas Ray-GPIS retains the reconstruction gap. In a
visited-but-unfused test, missing-patch Recall@5 improves from 0.2827 to 0.5066
(+79.2\%). Thus, pose coverage is effective under ideal sensing, while
reconstruction-conditioned planning is more robust when viewpoint coverage
and acquired geometry diverge.


\textbf{Q3 (AC, R2, R1): Is the improved mesh useful downstream?}
Yes. We run a sequential \textbf{quasi-static place-and-regrasp simulation} on ten metric
YCB objects. Each method has three paired reconstructions and ten execution
perturbations/mesh (30 trials/object). The task planner sees only the
reconstructed mesh, selecting a stable support pose and top-down antipodal
regrasp; perturbed execution is checked on GT geometry. Active methods use the
full five-action loop (6\,s/action, 15 FPS) with hand occlusion,
under-rotation, stall, and slip, and the same GPU-NKSR backend. TRELLIS.2
(released 1024-cascade geometry path) and SPAR3D receive a clean single image
and favorable oracle isotropic alignment.

Table~\ref{tab:downstream} shows that Ray-GPIS raises
placement/regrasp/joint success from 39.7/20.0/6.7\% with single RGB-D and
50.0/44.0/24.0\% with the fixed schedule to 57.7/77.7/45.0\%. A GT-mesh
oracle reaches 100/99/99\%, confirming task reachability. On six prior real
AURORA meshes evaluated against scanner GT, the rates are 66.7/72.2/41.7\%.

\begin{table}[t]
\centering
\caption{Controlled planner comparison and representative robustness tests.
Part (a) uses GT pose and exact RGB--depth registration; time is seconds per
planning step. Part (b) reports failure-specific metrics.}
\label{tab:main}
\vspace{-0.25cm}
\resizebox{\linewidth}{!}{%
\begin{tabular}{lccccc}
\toprule
\multicolumn{6}{l}{\textbf{(a) Ideal planner isolation: 120 paired episodes}} \\
Method & F-AUC$\uparrow$ & Final F$\uparrow$ & Corr.$\uparrow$ & Regret$\downarrow$ & Time$\downarrow$ \\
\midrule
Fixed schedule
  & $0.7948{\pm}0.0061$ & $0.9401{\pm}0.0046$ & --
  & $0.0806{\pm}0.0025$ & $0.000$ \\
Adapted ActNeRF~\cite{actnerf}
  & $0.8665{\pm}0.0062$ & $0.9651{\pm}0.0050$
  & $0.1619{\pm}0.0337$ & $0.0317{\pm}0.0029$
  & $2.238{\pm}0.018$ \\
Adapted PB-NBV~\cite{pbnbv}
  & $0.7475{\pm}0.0193$ & $0.8496{\pm}0.0234$
  & $-0.3563{\pm}0.0731$ & $0.1076{\pm}0.0124$
  & $0.391{\pm}0.024$ \\
Pose-Novelty
  & $\mathbf{0.9090{\pm}0.0043}$ & $\mathbf{0.9945{\pm}0.0014}$
  & $\mathbf{0.6186{\pm}0.0294}$ & $\mathbf{0.0025{\pm}0.0006}$
  & $2.49{\times}10^{-4}$ \\
\textbf{Full Ray-GPIS}
  & $0.8987{\pm}0.0091$ & $0.9814{\pm}0.0100$
  & $0.4722{\pm}0.0667$ & $0.0088{\pm}0.0051$
  & $0.261{\pm}0.024$ \\
\bottomrule
\end{tabular}%
}

\vspace{2pt}

\resizebox{\linewidth}{!}{%
\begin{tabular}{lllccc}
\toprule
\multicolumn{6}{l}{\textbf{(b) Targeted sensing-failure robustness tests}} \\
Stress & Comparator & Primary metric & Comp. & Full & Full benefit \\
\midrule
\multicolumn{6}{l}{\emph{Reconstruction-conditioned vs. pose-only planning}} \\
Visited but unfused depth (70\%) & Pose-Novelty
  & missing-patch Recall@5$\uparrow$ & 0.2827 & \textbf{0.5066}
  & \textbf{$+0.2239$ ($+79.2\%$)$^{\dagger}$} \\
\addlinespace[1pt]
\multicolumn{6}{l}{\emph{Component robustness ablations}} \\
Palm hole & w/o miss-ray interpolation & unseen Recall AUC$\uparrow$
  & 0.5498 & \textbf{0.5715}
  & \textbf{$+0.0217$ ($+4.0\%$)$^{\dagger}$} \\
Pose jitter ($6^\circ$/3\,mm) & w/o receptive-field integration
  & score-map correlation$\uparrow$ & 0.2724 & \textbf{0.5028}
  & \textbf{$+0.2304$ ($+84.6\%$)$^{\dagger}$} \\
\bottomrule
\end{tabular}%
}
\vspace{-2pt}
{\footnotesize Part (b) uses failure-specific metrics that are not
cross-row comparable. $^{\dagger}$Paired comparison remains significant after
Holm correction ($p<0.05$).}
\vspace{-0.15cm}
\end{table}

\begin{table}[t]
\centering
\caption{Sequential place-and-regrasp success (\%). Regrasp evaluates stage-2
feasibility in the intended settled pose; joint requires both stages.}
\label{tab:downstream}
\vspace{-0.22cm}
\resizebox{\linewidth}{!}{%
\begin{tabular}{lccc}
\toprule
Reconstruction source & Placement$\uparrow$ & Regrasp$\uparrow$ & Joint$\uparrow$ \\
\midrule
Single RGB-D & 39.7 & 20.0 & 6.7 \\
SPAR3D (oracle-align) & 30.0 & 15.0 & 5.0 \\
TRELLIS.2 (oracle-align) & 19.7 & 7.7 & 0.0 \\
\midrule
Fixed schedule & 50.0 & 44.0 & 24.0 \\
Adapted PB-NBV & 43.3 & 40.0 & 23.0 \\
Adapted ActNeRF & 56.0 & 72.0 & 41.7 \\
Pose-Novelty & 56.7 & \textbf{79.3} & 44.7 \\
\textbf{Full Ray-GPIS} & \textbf{57.7} & 77.7 & \textbf{45.0} \\
\midrule
GT mesh (oracle) & 100.0 & 99.0 & 99.0 \\
\bottomrule
\end{tabular}}
\vspace{-0.25cm}
\end{table}

\section*{2. Ablations, Robustness, Diversity \& Runtime}

\textbf{Q4 (R1, R3): Can the key Ray-GPIS components be ablated?}
Yes. We ablate uncertainty estimation, novelty weighting, miss-ray
interpolation, and receptive-field integration. On the clean benchmark,
removing uncertainty (Novelty Only) or novelty weighting (Uncertainty Only)
reduces F-AUC from 0.8987 to 0.8949 and 0.8953, respectively. Because the
remaining components target specific sensing failures, Table~\ref{tab:main}(b)
isolates those failure modes.
With a contiguous palm-induced hole, removing miss-ray interpolation discards
directions with no current surface return; Full retains these candidates and
raises unseen-surface Recall AUC from 0.5498 to 0.5715. Under
$6^\circ$/3\,mm pose jitter, removing receptive-field integration produces
fragmented pointwise scores, whereas Full raises score-map correlation from
0.2724 to 0.5028.

\textbf{Q5 (R1, AC): How robust is AURORA to 6D pose-tracking errors?}
AURORA is robust to mild-to-moderate errors but degrades under severe drift. We
perturb the tracked poses used by fusion and planning. As rotational/
translational error increases from $1^\circ$/1\,mm to $2^\circ$/2\,mm and
$10^\circ$/10\,mm, F@5 changes from 99.7\% to 98.8\% and 90.3\%, while
symmetric distance increases from 0.84 to 1.11 and 2.31\,mm.
Table~\ref{tab:main}(b) further shows that
receptive-field integration stabilizes the planner's score map under pose
perturbations, explaining this robustness.

\textbf{Q6 (R3, AC): Does the method generalize to more challenging
irregular/concave objects?} Yes. We add a concave Bowl and a Thin-Irregular
object (eight objects total), each evaluated from 15 poses. Their F-AUC/final
F@5 are $0.6402{\pm}0.0534$/$0.8107{\pm}0.0695$ and
$0.7046{\pm}0.0306$/$0.9274{\pm}0.0219$, respectively.

\textbf{Q7 (R3, R1, AC): What is the per-module runtime, and why replan
only every 6\,s?} We replan every 6\,s because this is the execution horizon of
one tactile rotation primitive, not because planning takes 6\,s. Segmentation
and end-to-end perception/BundleTrack take
$14.5{\pm}0.2$ and $242.9{\pm}47.9$\,ms/frame (the latter includes the former),
while fusion, GP update, candidate scoring, and NBV-to-action mapping take
$1,649.6{\pm}194.0$, $261{\pm}24$, $0.072{\pm}0.001$, and $2.42$\,ms/update,
respectively.

\section*{Other Responses}
\textbf{Q8 (R2): Why use in-hand rotation if Ray-GPIS is
controller-agnostic?} We use in-hand rotation to expose occluded surfaces while
retaining the grasp; controller agnosticism means that Ray-GPIS can also drive
another reorientation mechanism. Ray-GPIS selects \emph{which} under-observed
orientation to expose, while the in-hand controller determines \emph{how} to
reach it.

% \begin{figure}[t]
% \centering
% \vspace{-0.1cm}
% \fbox{\parbox{0.93\linewidth}{\centering\footnotesize\textit{[PLACEHOLDER FIGURE.} Top row: masked in-hand input, failed single-view
% NeRF geometry, AURORA reconstruction. Bottom row: a direction marked
% ``visited'' in pose space whose geometry stays incomplete and is retained as
% high-uncertainty by Ray-GPIS.\textit{]}\\
% \textit{Draw: 3 cropped RGB/depth/mesh triplets + 2 score-map spheres.}}}
% \vspace{-0.25cm}
% \caption{Visited-but-unfused: pose coverage $\neq$ acquired geometry.}
% \label{fig:qual}
% \vspace{-0.3cm}
% \end{figure}

{\footnotesize \bibliographystyle{plainnat}\bibliography{references}}

\end{document}
